\documentclass[ecta,nameyear,draft]{econsocart}
\usepackage{xr-hyper}
\usepackage{amsmath,amsfonts,amssymb,amsthm,graphicx,booktabs,array,longtable}
\usepackage{algorithm,algpseudocode,bm,enumitem}
\usepackage[T1]{fontenc}
\usepackage{mathpazo}
\usepackage{microtype}
\usepackage{needspace}
\RequirePackage[colorlinks,citecolor=blue,linkcolor=blue,urlcolor=blue]{hyperref}
\startlocaldefs
\newtheorem{theorem}{Theorem}
\renewcommand{\thetheorem}{S.\arabic{theorem}}
\renewcommand{\theequation}{S.\arabic{equation}}
\renewcommand{\thesection}{S.\arabic{section}}
\newtheorem{proposition}[theorem]{Proposition}
\newtheorem{lemma}[theorem]{Lemma}
\newtheorem{corollary}[theorem]{Corollary}
\newtheorem{assumption}{Assumption}
\theoremstyle{definition}
\newtheorem{definition}{Definition}
\theoremstyle{remark}
\newtheorem{remark}{Remark}
\newtheorem{example}{Example}
\newcommand{\E}{\mathbb E}
\newcommand{\R}{\mathbb R}
\newcommand{\HH}{\mathcal H}
\newcommand{\LL}{\mathcal L}
\newcommand{\Act}{\mathcal A}
\newcommand{\sg}{\operatorname{sg}}
\newcommand{\tr}{\operatorname{tr}}
\newcommand{\Lip}{\operatorname{Lip}}
\DeclareMathOperator*{\argmax}{arg\,max}
\DeclareMathOperator*{\argmin}{arg\,min}
\endlocaldefs

\let\NBOLongtable\longtable
\def\longtable{\Needspace{7\baselineskip}\NBOLongtable}

\setlength{\emergencystretch}{2em}
\externaldocument{revisions/2026-10-07-r39/build/supp_from_historical_article}[historical_article.pdf]
\externaldocument{revisions/2026-10-07-r39/build/supp_from_historical_supplement}[historical_supplement.pdf]
\externaldocument{revisions/2026-10-07-r39/build/supp_from_applications}[applications.pdf]
\externaldocument{revisions/2026-10-07-r39/build/supp_from_ECTA}[ECTA.pdf]
\externaldocument{revisions/2026-10-07-r39/build/supp_from_response}[response.pdf]

\makeatletter
\def\copyright@text{Prepared for referee review}
\makeatother
\makeatletter\def\copyright@text{Prepared for referee review}\makeatother
\begin{document}
\begin{frontmatter}
\title{Technical Supplement to Neural Bellman Operators}
\runtitle{Neural Bellman Operators}
\begin{aug}
\author[id=au1,addressref={add1}]{\fnms{Qian}~\snm{QI}}
\address[id=add1]{Peking University}
\end{aug}
\end{frontmatter}
\section{Decision Accuracy, Transport, and Stopping: Proofs}\label{app:r19proofs}
This section gives the complete arguments for the new economic error chain.
The inherited verification, Bellman-bridge, and diffusion proofs are retained
in the preceding sections. The constants below pertain to the new stored
finite-law experiment unless a more general hypothesis is stated.

\subsection{Proof of Proposition~\ref{prop:r19decision}}
First suppose an optimal action $a^*$ exists. The approximate maximizing
property implies
\[
 r_\theta(s,a^*)+\widehat S_0(T_\theta(s,a^*))
 -r_\theta(s,\widehat a)-\widehat S_0(T_\theta(s,\widehat a))
 \leq\delta_\theta(s).
\]
Substituting $S_\theta=\widehat S_0-e+\Delta_\theta$ gives
\begin{align*}
 Q_\theta(s,a^*)-Q_\theta(s,\widehat a)
 \leq {}&\delta_\theta(s)
       +e(T_\theta(s,\widehat a))-e(T_\theta(s,a^*))\\
       &+\Delta_\theta(T_\theta(s,a^*))
                -\Delta_\theta(T_\theta(s,\widehat a)).
\end{align*}
Each difference is at most its action-set oscillation. For an unattained
finite supremum, take an $\epsilon$-maximizer and let $\epsilon$ decrease
to zero. This proves \eqref{eq:r19osc} without a differentiability assumption.

For finite actions let $v_j=e(T_\theta(s,a_j))-\bar e_\pi(s)$.
In the weighted inner product $\langle x,y\rangle_\pi=\sum_j\pi_jx_jy_j$,
the representer of $v_i-v_j$ has entries $1/\pi_i$ at $i$, $-1/\pi_j$
at $j$, and zero elsewhere. Its squared norm is $1/\pi_i+1/\pi_j$.
Cauchy--Schwarz gives
\[
 (v_i-v_j)^2\leq(1/\pi_i+1/\pi_j)\sum_k\pi_kv_k^2.
\]
Maximizing over the finite pair set proves
$\operatorname{osc}e\leq\sqrt{C_\pi R_\pi}$. The transport difference is
at most $2b_\theta$, proving \eqref{eq:r19riskloss}. For random fitted
objects the inequality is pathwise. With deterministic $C_\pi$, for example,
$\E\sqrt{C_\pi R_\pi}\leq\sqrt{C_\pi\E R_\pi}$ whenever the expectation
is finite. Repetition of the same random seed does not create a new draw
for this expectation. Finally, if every nonoptimal action loses at least
$g$ and the displayed bound is below $g$, a nonoptimal implemented action
would contradict the bound.\hfill$\square$

\subsection{Proof of Theorem~\ref{thm:r19transport}}
Couple both economies by the same innovation sequence and the same initial
postdecision state. If $\|X_k^\theta-X_k^0\|\leq d_k$, then the triangle
inequality and \eqref{eq:r19transportassumptions} give
\[
 \|X_{k+1}^\theta-X_{k+1}^0\|
 \leq\epsilon_{\theta,k}+L_kd_k=d_{k+1}.
\]
Induction starts at $d_0=0$. At each reward date,
$|g_{\theta,k}(X_k^\theta)-g_{0,k}(X_k^0)|
 \leq\eta_{\theta,k}+H_kd_k$.
Multiply by the nonnegative weights, sum including the terminal reward,
and take expectations. The bound is uniform because each assumed bound is
uniform on the invariant domain. No interchange of a supremum over future
controls with a fixed-policy expectation is used.

For recursive evaluations write
$D_k=\|V_{\theta,k}-V_{0,k}\|_\infty$. Add and subtract
$\mathcal G_{\theta,k}(V_{0,k+1})$ to obtain
$D_k\leq\nu_{\theta,k}+\beta_kD_{k+1}$.
Backward substitution from $D_K\leq b_K$ gives
\eqref{eq:r19recursive}. Both arguments require that every coupled state
and every recursively evaluated utility belong to the domains on which the
bounds were verified. The theorem makes no extension past these domains.
If a domain is localized, an exit contribution must be bounded separately
before using the result as a global guarantee.\hfill$\square$

\subsection{Proof of Proposition~\ref{prop:r19stopping}}
Strong concavity gives, for any feasible $b$,
\[
 Q(b)-Q(a)\leq Q'(a)(b-a)-\kappa(b-a)^2/2.
\]
In the interior, maximize this upper bound over all real $b-a$ to obtain
$|Q'(a)|^2/(2\kappa)$. At the lower endpoint restrict the maximization to $b-a\geq0$.
Its value is
\[
 \frac{(\max\{Q'(a),0\})^2}{2\kappa}.
\]
At the upper endpoint the analogous residual is $\max\{-Q'(a),0\}$. Enlarging a bounded
feasible half-interval to the half-line is conservative. This proves
\eqref{eq:r19residual}.

An outward derivative interval bounds the applicable residual, and a positive
lower curvature bound enlarges the resulting nonnegative fraction. Thus each
task certificate holds deterministically at every attempted stage. A first
successful average certificate implies the average economic target; previous
failures do not affect validity. If checks are stochastic, let $B_j$ denote
the event that attempted check $j$ falsely certifies its candidate. The
conditional hypothesis gives $\Pr(B_j)\leq\E\alpha_j$; an unattempted
check contributes zero. The union bound yields
$\Pr(\bigcup_jB_j)\leq\E\sum_j\alpha_j\leq\alpha$. This proof requires
conditional validity for the data actually used after past decisions.
A fixed-time interval repeatedly queried without such a property does not
meet the hypothesis.

The work identity follows by summing the disjoint measured operations through
the stop, including unsuccessful fits and tests. Under the stated exact
linear cost model, subtraction gives
$W_N(q)-W_R(q)=F_N-F_R-q(c_R-c_N)$, establishing the strict inequality
in \eqref{eq:r19break}. Integer query counts use that strict inequality;
a tie at an integer threshold is not a strict saving. Measured costs in the
experiment are used directly and are not assumed to satisfy a linear model.
\hfill$\square$

\subsection{Proof of Corollary~\ref{cor:r19amortization}}
Condition on the feature map and design, which are independent of the fitting
noises. Least squares gives
$\widehat\beta-\beta=G_n^{-1}\Phi^{\mathsf T}\xi$.
For a fixed queried argument $x$, the prediction error is a linear combination
of the independent noises. Its sub-Gaussian variance parameter is at most
\[
 \sigma^2\|\Phi G_n^{-1}\phi(x)\|^2
 =\sigma^2\phi(x)^{\mathsf T}G_n^{-1}\phi(x)
 \leq\sigma^2\ell/n.
\]
The two-sided sub-Gaussian bound and a union bound over the $K$ distinct
arguments therefore give
\[
 \Pr\left(\max_x|\widehat S(x)-S(x)|>
  \sigma\sqrt{\frac{2\ell}{n}\log\frac{4K}{\alpha}}\right)
 \leq\alpha/2.
\]
Uniformly on that event, the oscillation in every feasible finite action
menu is at most twice the displayed bound. Proposition~\ref{prop:r19decision}
with zero transport and \eqref{eq:r19learnbudget} gives loss at most
$\varepsilon$ at every query. Full column rank and the leverage bound are
separate hypotheses and must actually hold at the selected integer $n$.
They are not consequences of the scalar sample-size inequality.

For the direct method, each prediction average has sub-Gaussian variance
parameter at most $\sigma^2/t$. Dependence between predictions caused by
common innovations across different arguments does not affect the union
bound; independence is required only for the observations within an average,
or an equivalent sub-Gaussian bound for that average. The same calculation
with $\ell/n$ replaced by $1/t$ proves the second guarantee. A union bound
over the two methods gives joint coverage $1-\alpha$, without assuming
independence between them. Reusing a cached prediction at exactly the same
continuation identity and argument incurs no additional simulation count.
This is why $K$, rather than the number of repeated requests, enters the
work comparison.

Subtract the complete specified costs to obtain
\eqref{eq:r19explicitadvantage}. The fixed term $F(n)$ includes dictionary
construction, all labels, the least-squares solve and its validation; it is
not just a neural update clock. Common selection and verification work
cancels only under the explicit equal-cost hypothesis. Otherwise its
difference remains. These are sufficient budgets and actual costs for the
declared procedures, not statistical lower bounds on arbitrary algorithms.
If $|\E\widehat S(x)-S(x)|\leq B$ uniformly under a misspecified dictionary,
the prediction bound acquires $B$; if future continuation changes by at most
$b$, the decision bound acquires a further $2b$. Thus the same argument
uses margin $\varepsilon-\delta-2B-2b>0$, provided the stated bias bound is
verified. A pointwise approximation error of an arbitrary basis expansion
is not automatically a bound on the bias of its fitted least-squares
projection.\hfill$\square$

\subsection{Uniform curvature in the stored capital economy}
\label{app:r19curvature}
Let $M=(1-\mu)I+\mu\bm1\bm1^{\mathsf T}/d$. It is row stochastic and
has infinity operator norm one. For $f(y)=\lambda\tanh(My)$, the inequalities
$|\tanh'|\leq1$ and $|\tanh''|\leq2$ imply, along a scalar path with
$\|y'\|_\infty\leq P$ and $\|y''\|_\infty\leq D$,
\begin{equation}
 \|(f(y))''\|_\infty\leq\lambda(D+2P^2).
 \label{eq:r19prodcurv}
\end{equation}
The postdecision map has first derivative $-h\bm1$ and second derivative
zero, so initialize $P=h$, $D=0$. Adding the first innovation does not
change these derivatives. For each of the three future transitions set
\begin{align*}
 B_k&=\lambda(D+2P^2),&
 D&\leftarrow D+hB_k,&P&\leftarrow(1+h\lambda)P.
\end{align*}
The order matters: $B_k$ uses the preceding $P,D$. The second derivative of
the future production reward is at most $W_kB_k$, since every stored $W_k$
is positive. All future terms involving the fixed withdrawal are constant
with respect to the current action.

For a vector $v$ write $\operatorname{range}(v)=\max_i v_i-\min_i v_i$.
At the first shocked postdecision state a valid range bound for each stored
path is
\[
 R_0=\operatorname{range}(y)+2h\lambda+
                       \operatorname{range}(z_0).
\]
Every subsequent transition satisfies
$R_k=R_{k-1}+2h\lambda+\operatorname{range}(z_k)$ as an upper bound.
The scalar drift and common withdrawal cancel from a coordinate range.
For the terminal payoff $g(y)=-\gamma\operatorname{Var}_d(y)/2$,
\[
 (g(y))''=-\gamma\operatorname{Var}_d(y')
 -\gamma\overline{(y-\overline y\bm1)y''}
 \leq\gamma R_3D_3.
\]
The first term is nonpositive; dropping it is an upper bound, not a sign
reversal. Combining these inequalities gives a uniform upper bound on the
continuation curvature,
\[
 S_0(x(y,a))''\leq
 M(y):=\sum_{k=1}^{3}W_kB_k+\gamma\E_{16}[R_3D_3].
\]
Since the current reward has second derivative
$-A_0(w/a^2+\eta)$, a valid lower bound for true strong concavity on the
whole interval is
\begin{equation}
 \kappa(y,\theta)=A_0\{w/\bar a^2+\eta\}-M(y).
 \label{eq:r19kappa}
\end{equation}
The verifier evaluates this expression with outward arithmetic and requires
a strictly positive lower endpoint for every task. A nonpositive result
would be a verification failure, not a reason to delete that task. The
published records retain every returned lower endpoint.

The true derivative is enclosed separately rather than inferred from the
surrogate. For each stored path initialize $p_0=-h\bm1$. At a future state,
\[
 f'(y)=\lambda\{1-\tanh^2(My)\}Mp,\qquad
 p\leftarrow p+h f'(y).
\]
Accumulate $W_k\overline{f'(y)}$ and the terminal derivative
$-\gamma\overline{(y-\overline y\bm1)p}$. Average over all sixteen atoms
and add the current derivative
$A_0(w/a-\eta a-\tau)-W_0$. Every operation is evaluated on enclosing
intervals. Dependency between occurrences of the same variable may widen
an interval, but cannot invalidate the enclosure.

\subsection{Arithmetic enclosure and its implementation}
The stored binary64 coefficients and innovation increments define the exact
finite economy. Floating-point evaluation of a Gaussian generator is not
being asserted to be exact sampling from a continuous Gaussian economy.
The generator is used to construct and freeze a finite array; the subsequent
scientific statement concerns that array with equal masses.

The arithmetic module encloses each binary operation by the adjacent
representable numbers toward minus and plus infinity. It checks finite
values, IEEE round-to-nearest behavior and gradual underflow. Products and
quotients take extrema over endpoint combinations, with divisors bounded
away from zero. Squaring includes zero when the argument interval contains
zero. Balanced reductions avoid interpreting a non-enclosing ordinary mean
as an interval sum. Induction on an expression's operation graph proves
inclusion for arithmetic expressions from enclosed inputs.

The nonlinear enclosure does not assume that a library's hyperbolic tangent
is correctly rounded. For $0\leq t\leq1/2$, the exponential is enclosed by
its degree-18 Taylor polynomial and the nonnegative tail bound
\[
 0\leq e^t-\sum_{j=0}^{18}t^j/j!
 \leq \frac{t^{19}/19!}{1-t/20}.
\]
All terms are enclosed using the same arithmetic. Exact binary scaling (checked on every replayed argument) reduces the
argument to this interval; repeated squaring and reciprocation recover
positive and negative arguments. The identity
$\tanh x=(e^{2x}-1)/(e^{2x}+1)$ and monotonicity provide endpoint enclosures.
For the auxiliary logarithm check, scaling to $[1,2)$ gives
$z=(x-1)/(x+1)\in[0,1/3)$ and
\[
 \log x=2\sum_{j=0}^{31}\frac{z^{2j+1}}{2j+1}+R,
 \qquad 0\leq R\leq\frac{2z^{65}}{65(1-z^2)}.
\]
The stored enclosing endpoints for $\log2$ is independently checked with
exact rational series bounds. Synthetic high-precision derivative tests
check the implemented recursion on separate small-dimensional cases.
These tests complement the analytical enclosure argument; they are not
substitutes for it or additional economic observations.

\subsection{Charge monotonicity and the economic tolerance}
For an interior maximizing action, differentiate
$Q_\theta'(y,a^*(\tau))=0$. Because the charge enters the derivative as
$-A_0\tau$, the implicit-function identity is
$Q''\,\partial a^*/\partial\tau-A_0=0$, proving
\eqref{eq:r19charge}. For the constrained case take $\tau_2>\tau_1$ and
respective optimal actions $a_2,a_1$. Adding their optimality inequalities
leaves
$A_0(\tau_2-\tau_1)(a_2-a_1)\leq0$, hence $a_2\leq a_1$.
Uniqueness follows from strict concavity.

An external grant in the utility account satisfies
$A_0w\log(1+g)=U$ exactly when $g=\exp(U/(A_0w))-1$. This is a
monotone normalization of a certified payoff loss. The published grant
table uses ordinary binary64 evaluation and is explicitly descriptive;
the primary rigorous statement is the outward payoff-loss bound. No claim
is made that increasing actual endogenous consumption by this amount leaves
capital, future rewards, or the reference policy unchanged.

\section{From Current Decisions to Policy Accuracy}\label{sec:r21composition}
A continuation certificate is useful for a dynamic economy only if the
certified object is the future of the policy under assessment. This section
states the conditions under which the local decision accounts compose.
They apply to additive and order-preserving recursive evaluation, while
keeping the utility domains and admissible deviations of the economic model.

At dates $t=0,\ldots,T-1$, let $\mathcal D_t$ be an invariant state domain
and $A_t(s)$ the full admissible current action set. Write
\[
 (T_t v)(s)=\sup_{a\in A_t(s)}G_t(s,a,v),\qquad
 (T_t^\pi v)(s)=G_t(s,\pi_t(s),v).
\]
Assume the dynamic programming principle on these domains, measurable
admissible policies, and finite values. Each $G_t$ is order preserving in its
continuation argument and is sup-norm Lipschitz, uniformly in $(s,a)$, with
constant $\beta_t\geq0$ on a common invariant value domain. Let
$V_T^*=V_T^\pi$ and define $V_t^*=T_tV_{t+1}^*$ and
$V_t^\pi=T_t^\pi V_{t+1}^\pi$. The relevant local advantage is
\begin{equation}
 g_t^\pi(s)=T_tV_{t+1}^\pi(s)-T_t^\pi V_{t+1}^\pi(s).
 \label{eq:r21advantage}
\end{equation}
It evaluates all current deviations against the \emph{same implemented
future policy}; it is neither a training residual nor an advantage computed
against an unrelated anchor continuation.

\begin{theorem}[Composition of economic certificates]\label{thm:r21composition}
Under the preceding assumptions, suppose that on a common event
\begin{equation}
 0\leq g_t^\pi(s)\leq\varepsilon_t
 \quad\text{for every }s\in\mathcal D_t,
 \qquad t=0,\ldots,T-1.
 \label{eq:r21uniform}
\end{equation}
Then, on that event, at every initial state,
\begin{equation}
 0\leq V_0^*-V_0^\pi\leq
 \sum_{t=0}^{T-1}w_t\varepsilon_t,
 \qquad w_t=\prod_{j=0}^{t-1}\beta_j.
 \label{eq:r21composition}
\end{equation}
For additive Markov evaluation
$G_t(s,a,v)=r_t(s,a)+\beta_tP_t^av(s)$, a state-dependent refinement is
$V_t^*-V_t^\pi\leq B_t$, where
\begin{equation}
 B_T=0,\qquad
 B_t(s)=u_t(s)+\beta_t\sup_{a\in A_t(s)}P_t^aB_{t+1}(s)
 \label{eq:r21envelope}
\end{equation}
whenever $u_t(s)\geq g_t^\pi(s)$ and the displayed quantities are finite.
\end{theorem}

Section~\ref{app:r21proofs} proves the result. Centered continuation-error
and transport bounds from Proposition~\ref{prop:r19decision} can supply
\eqref{eq:r21uniform} when established uniformly on the required state and
action domains. Alternatively, an interval certificate for the true
one-step objective can supply it directly. In either case, the continuation
must be $V_{t+1}^\pi$, or its discrepancy must be charged. The verification
may be deterministic. With random certificates, all domain and date bounds
must hold on one event, for example by a prespecified summable conditional
error allocation. Selecting a stopping time does not remove this requirement.

This theorem identifies a positive route from reusable policy evaluation to
full policy accuracy. It also identifies precisely what the computational
evidence below supplies. A mean certificate on a finite state catalogue
establishes the declared current-decision experiment; it does not establish
\eqref{eq:r21uniform} outside that catalogue. Sampling only states visited by
$\pi$ is insufficient because a profitable deviation may reach other states.
The state envelope in \eqref{eq:r21envelope} deliberately takes a supremum
over admissible transitions, not only the transitions of the fitted policy.

\paragraph{Relation to the continuous economy and strategic applications.}
Let $\Pi_h$ be a specified implementable policy class in the original economy.
If its class-approximation gap is at most $\eta_{\rm class}$, and
$|J(\pi)-J_h(\pi)|\leq\eta_{\rm value}$ uniformly for $\pi\in\Pi_h$, then a
policy satisfying \eqref{eq:r21composition} in the approximating model has
\begin{equation}
 \sup_{\pi}J(\pi)-J(\widehat\pi)
 \leq\eta_{\rm class}+2\eta_{\rm value}
       +\sum_t w_t\varepsilon_t+\eta_{\rm impl},
 \label{eq:r21bridge}
\end{equation}
where $\eta_{\rm impl}$ bounds any additional payoff loss in implementation.
The class and value-transfer allowances are separate mathematical inputs,
not numbers inferred from the finite experiments. For recursive utility,
order and Lipschitz bounds are imposed on the actual common utility domain;
actionwise centering is applied to the evaluated action payoff, not through
an arbitrary nonlinear aggregator. For a dynamic game, apply the same
account to each player's full deviation problem against the same opponents'
profile. Bounds $\epsilon_i$ for all players imply an $\epsilon_i$-best-response
profile under the specified equilibrium concept. A single local deviation
or a single player's calculation does not supply those premises. The original
information, boundary, transversality, and deviation requirements remain in
the Economic Applications companion.

\subsection{Allocating computation across dates}
The policy target suggests how to distribute a verified approximation budget.
Suppose the local account at date $t$ is $b_t+e_t$, where $b_t$ contains
fixed optimization, representation, and transport allowances. Suppose further
that a declared procedure attains approximation allowance $e_t>0$ using
$n_t\geq A_t/e_t^2$ fitting observations at charged marginal cost $c_t>0$.
Here $A_t>0$ includes its confidence allocation and design constants. The
full-rank and coverage conditions underlying such a rate must hold at the
selected budgets; Corollary~\ref{cor:r19amortization} is one finite-query
example, not a license to assume uniform coverage of an infinite state space.

\begin{proposition}[Allocation for a policy target]\label{prop:r21allocation}
Assume $w_t>0$ and let
$E=\varepsilon-\sum_t w_tb_t>0$, $a_t=c_tA_t$, and
$H=\sum_ta_t^{1/3}w_t^{2/3}$. Among positive allowances satisfying
$\sum_tw_te_t\leq E$, the minimum continuous fitting bill is
\begin{equation}
 \min\sum_t\frac{a_t}{e_t^2}=\frac{H^3}{E^2},\qquad
 e_t=\frac{E}{H}\left(\frac{a_t}{w_t}\right)^{1/3}.
 \label{eq:r21allocation}
\end{equation}
Rounding each observation count up adds at most $\sum_tc_t$ to this bill.
All fixed construction, reusable-cache, query, and verification costs must
be added. The statement minimizes the displayed sufficient-cost account,
not the work of all possible numerical methods.
\end{proposition}

The propagation and allocation arguments use standard dynamic-programming
stability and convex optimization. Their role here is to connect the paper's
centered, transported, and true-objective certificates to a specified dynamic
economic target. They do not replace approximation or coverage assumptions.
Finite-time fitted-value analysis, including the distinction between sampling
norms and policy loss, provides an important antecedent
\citep{MunosSzepesvari2008}. The empirical contribution below remains a
measured comparison of complete procedures, rather than a claim of a new
principle of optimality.

\section{Proofs for Policy Composition and Economic Contrasts}\label{app:r21proofs}
\begin{proof}[Proof of Theorem~\ref{thm:r21composition}]
Backward induction and order preservation give $V_t^*\geq V_t^\pi$.
Suppose $0\leq V_{t+1}^*-V_{t+1}^\pi\leq b$ on $\mathcal D_{t+1}$.
The elementary supremum inequality and the uniform Lipschitz condition give
\[
 0\leq T_tV_{t+1}^*-T_tV_{t+1}^\pi\leq\beta_tb.
\]
Adding $T_tV_{t+1}^\pi-V_t^\pi=g_t^\pi$ shows that the sup-norm error is at
most $\varepsilon_t+\beta_tb$. Iteration from the common terminal value
proves \eqref{eq:r21composition}, including the case of a zero Lipschitz
constant. No assumption that $\pi$ itself attains the optimal action is used.
For the state envelope, suppose by induction that
$V_{t+1}^*-V_{t+1}^\pi\leq B_{t+1}$. Then
\begin{align*}
 V_t^*(s)-V_t^\pi(s)
 &\leq \sup_a\{r_t(s,a)+\beta_tP_t^aV_{t+1}^\pi(s)
                  +\beta_tP_t^aB_{t+1}(s)\}-V_t^\pi(s)\\
 &\leq g_t^\pi(s)+\beta_t\sup_aP_t^aB_{t+1}(s)
 \leq B_t(s).
\end{align*}
Invariance ensures that every expectation is taken on a domain on which the
induction applies. These are pathwise deductions on the common certificate
event, so no additional probability statement is introduced by the proof.
\end{proof}

\paragraph{Why catalogue averages cannot be substituted.}
Consider a two-date economy with zero first-date rewards. A deviation at the
initial state can reach a state with a unit second-date gain that the
implemented policy never visits. All local advantages at the implemented
policy's visited states can be zero while the initial optimality gap equals
one. Thus a bound only on that occupancy distribution cannot be substituted
for \eqref{eq:r21uniform}. The same example shows why changing the supremum
in \eqref{eq:r21envelope} to the transition of $\pi$ is generally invalid.
Weighted sampling-norm results require their own distribution-shift or
coverage assumptions. We make no such inference from the finite query banks.

\begin{proof}[Proof of the bridge \eqref{eq:r21bridge}]
For the returned approximating policy $\pi_h\in\Pi_h$, write
\begin{align*}
 \sup_\pi J(\pi)-J(\widehat\pi)
 &\leq\eta_{\rm class}
       +\sup_{\pi\in\Pi_h}J(\pi)-J(\pi_h)+\eta_{\rm impl}\\
 &\leq\eta_{\rm class}+2\eta_{\rm value}
       +\sup_{\pi\in\Pi_h}J_h(\pi)-J_h(\pi_h)+\eta_{\rm impl}.
\end{align*}
Apply Theorem~\ref{thm:r21composition} to the last difference with the
specified policy class and action information. If an enlarged approximating
action class is used, its optimum is an upper bound and the inequality
continues to hold. The conclusion for fixed opponents follows by taking this
same supremum over each player's admissible deviations. A simultaneous
profile of these inequalities is precisely the stated approximate
best-response condition; no existence or fixed-point claim is inferred.
\end{proof}

\begin{proof}[Proof of Proposition~\ref{prop:r21allocation}]
For any positive feasible allowances, H\"older's inequality implies
\[
 H=\sum_t\left(\frac{a_t}{e_t^2}\right)^{1/3}
                    (w_te_t)^{2/3}
 \leq\left(\sum_t\frac{a_t}{e_t^2}\right)^{1/3}
                    \left(\sum_tw_te_t\right)^{2/3}.
\]
Hence the objective is at least $H^3/E^2$. The displayed allowances satisfy
$\sum_tw_te_t=E$ and attain equality. Since
$0\leq\lceil A_t/e_t^2\rceil-A_t/e_t^2<1$, their integer observation cost
exceeds the continuous bill by less than $\sum_tc_t$. At any date whose
statistical design requires a minimum sample size, that requirement must
also hold at the rounded budget; otherwise the design must be charged and
the constrained allocation recomputed. Dates with zero propagation weight
are omitted from this positive-weight allocation problem.
\end{proof}

\subsection{Optimal-action intervals and comparative statics}
The economic bands in Table~\ref{tab:r21nbocontrasts} use more information
than a mean value-loss certificate. Let $Q$ be continuously differentiable
and $\kappa$-strongly concave on $[\ell,u]$, with unique maximizer $a^*$.
For feasible $a$, let $r(a)$ be the directional residual in
Proposition~\ref{prop:r19stopping}. Then
\begin{equation}
 |a-a^*|\leq r(a)/\kappa.
 \label{eq:r21actionradius}
\end{equation}
To see this for $a<a^*$, the optimality condition gives $Q'(a^*)\geq0$
when $a^*=u$, and gives zero at an interior optimum. Strong concavity then
implies $Q'(a)\geq\kappa(a^*-a)$. For $a>a^*$, the reverse argument gives
$-Q'(a)\geq\kappa(a-a^*)$. At an endpoint the positive-part residual gives
exactly the appropriate one-sided inequality. The case $a=a^*$ is immediate.

Use the interval upper bound $\overline r_i$ and the strictly positive
curvature lower bound $\underline\kappa_i$ for each task. Outward arithmetic
and intersection with feasibility give
\[
 I_i=\left[a_i-\frac{\overline r_i}{\underline\kappa_i},
           a_i+\frac{\overline r_i}{\underline\kappa_i}\right]
                   \cap[\ell_i,u_i],\qquad a_i^*\in I_i.
\]
For the same tasks under a changed future $\theta$ and the anchor, the mean
optimal response belongs to $q^{-1}\sum_i(I_{\theta,i}-I_{0,i})$.
No independent-noise assumption is needed: these are deterministic
inclusions for fixed finite laws. We use the final NBO actions, not the
enumeration actions, to obtain the NBO response bands. The envelope over
three fitting streams contains every corresponding inclusion; it does not
turn three fixed streams into a population confidence interval. Printed
endpoints are rounded outwards. A negative upper endpoint establishes a
negative mean optimal response in that declared economy.

\subsection{Source and arithmetic audit}
The R21 audit reads every original R20 service record, verifies its stored
SHA256 digest, reestablishes the inherited source identities, regenerates
only the fixed task arguments, and recomputes the certificates at the saved
actions. It does not retrain any predictor or replace any clock. All
984 attempted certificates and 120 additional zero-charge certificates are
compared field for field with their original values. The attempted checks
cover 319,062 task actions. The audit also verifies feasibility, stopping
order, recorded success flags, every cost decomposition, and the exact
binary scaling used in the interval exponential implementation. This is an
arithmetic and source audit of the specified finite model, not an empirical
assessment of external economic calibration.

\section{Full-Policy Verification Details}\label{sec:r23proofs}
This section supplies the arithmetic and implementation accounts for
Section~\ref{sec:r23policy}. All numerical coefficients are the exact values
of their stored binary64 representations. The stochastic law remains
Gaussian with those coefficients; no finite shock catalogue is substituted.

\subsection{The quadratic local advantage}
For a fixed supplied future policy, substitute
$U_{t+1}(z)=z'P_{t+1}z+c_{t+1}$ into its one-step cost. The action-dependent
terms are
\[
 u'H_tu+2\beta u'B_t'P_{t+1}A_tx.
\]
Since $H_t\succ0$, the minimum over the entire space $\R^d$ occurs at
$u_t^g=-H_t^{-1}\beta B_t'P_{t+1}A_tx$. Evaluating at $-K_tx$ and subtracting
this minimum yields
\[
 (-K_tx-u_t^g)'H_t(-K_tx-u_t^g)=x'D_t'H_t^{-1}D_tx.
\]
There is no approximation in this identity and no reference to the optimal
future policy. The possibly loose Frobenius bound in
\eqref{eq:r23eta} is chosen to make the implementation independent of an
uncertified spectral decomposition. A failed bound need not indicate that
the actual policy loss exceeds the tolerance.

\subsection{Outward matrix balls}
A matrix ball $(C,R)$ denotes every matrix $M$ with
$|M-C|\leq R$ entrywise, where $R\geq0$. Exact input balls have zero radius.
For conformable balls $(C,R)$ and $(D,S)$, let $k$ be the inner dimension.
The exact product is enclosed about the computed product $\operatorname{fl}(CD)$
by the radius
\begin{equation}
 \gamma_{2k}|C||D|+|C|S+R|D|+RS+2k\tau\mathbf1,
 \label{eq:r23ballproduct}
\end{equation}
with every nonnegative calculation rounded upward. Here
$\gamma_n=nu/(1-nu)$, $u$ is the binary64 unit roundoff, and $\tau$ is the
smallest positive normal binary64 number. The underflow allowance deliberately
exceeds the corresponding subnormal absolute-error allowance. The factor
$2k$ permits separate multiplication and addition rounding in classical
binary64 products. Fused operations only reduce this particular account.
Overflow and nonfinite values are rejected rather than certified.

To calculate a nonnegative product upper bound, the implementation divides
the computed product by a downward bound for $1-\gamma_{2k}$, adds the
underflow allowance, and rounds upward. Addition, scalar multiplication,
traces, squared Frobenius norms, and maximum absolute row sums have analogous
outward bounds. These operations enclose the actual recursion, not merely
a rounded residual evaluated without error bars. Positivity of $P_t$ follows
from the exact evaluation equation and positive stage costs, not from an
unverified numerical eigenvalue. In the executed model the positive diagonal
entries of $Q_t,R_t$ give the coercivity lower bounds directly. Covariance
validity is checked by symmetry and strict diagonal dominance.

The tests compare matrix operations with rational arithmetic on the exact
binary input values, including signed products and cancellation. Separate
tests compare policy losses with a separately computed optimal comparator,
check the coercive bound at states far outside an ordinary training range,
perturb policies within the declared execution class, and verify that the
certificate still runs when the Riccati comparator is disabled. These tests
complement, rather than replace, the enclosure argument.

\subsection{A robust action-implementation envelope}
Suppose the implemented action is $-K_tx+e_t$, with
$\|e_t\|_2\leq e\|x\|_2$, possibly depending on the past and current state.
Put $\overline P_T=Q_T$ and $\overline c_T=0$. For $F_t=A_t-B_tK_t$, set
\begin{align}
 h_t={}&2\|R_tK_t\|_F e+\|R_t\|_\infty e^2\nonumber\\
 &+\beta\{2\|\overline P_{t+1}F_t\|_F\|B_t\|_F e
             +\|\overline P_{t+1}\|_\infty\|B_t\|_F^2e^2\},
 \label{eq:r23implh}\\
 \overline P_t={}&Q_t+K_t'R_tK_t+
                    \beta F_t'\overline P_{t+1}F_t+h_tI,\nonumber\\
 \overline c_t={}&\beta\{\overline c_{t+1}+\tr(\overline P_{t+1}\Sigma)\}.
 \label{eq:r23implrec}
\end{align}
All displayed coefficients are replaced by their outward upper enclosures
in the numerical calculation.

\begin{proposition}[Implementation allowance]\label{prop:r23implementation}
For every execution satisfying the relative action-error bound, its remaining
cost is at most $x'\overline P_tx+\overline c_t$. On the initial ball
$\|x\|_2^2\leq b$, its excess over the ideal stored policy is at most
\begin{equation}
 \varepsilon_{\rm impl}
 =\max\{0,b\|\overline P_0-P_0\|_\infty+
                         \overline c_0-c_0\}.
 \label{eq:r23implallow}
\end{equation}
\end{proposition}
\begin{proof}
Expanding the current action cost around $-K_tx$ bounds its increase by
$\{2\|R_tK_t\|_F e+\|R_t\|_\infty e^2\}\|x\|_2^2$.
Expanding the next-state quadratic envelope around $F_tx$ bounds its
increase by the second line of \eqref{eq:r23implh}, times $\|x\|_2^2$.
The mean-zero innovation contributes its trace term and has no cross term
with a nonanticipating action. Backward induction proves the value envelope.
Subtract the exact stored-policy quadratic value, use symmetry and the
maximum-row-sum bound, and maximize over the initial ball. This proves
\eqref{eq:r23implallow} even when the execution errors are history dependent.
\end{proof}

The error budget is explicit. Certification of arbitrary machine executions
without an error contract would be a different claim. The experiment neither
removes this term nor infers it from agreement on sampled states.

\subsection{Reoptimized comparative statics}
Let $K^{(r)}$ be a certified complete policy in regime $r$, and let its
reported total loss bound on the initial ball be $b_r$. At a fixed initial
state, the first-action objective with the \emph{optimal} future is a strictly
convex quadratic with Hessian at least $2R_0$. The cost difference between
using the candidate's first action and the optimal first action, both followed
by the optimal future, is no larger than the candidate's full policy loss.
Thus, with $r_0\leq\lambda_{\min}(R_0)$,
\begin{equation}
 \|u_0^{(r)}-u_0^{*,(r)}\|_2\leq\sqrt{b_r/r_0}.
 \label{eq:r23actionradius}
\end{equation}
At $x_0=-\mathbf1$, define mean investment
$m_r=d^{-1}\mathbf1'K_0^{(r)}\mathbf1$. Its distance from the optimal mean
is at most $\sqrt{b_r/(r_0d)}$. The reported radius also includes the
relative-execution and summation allowances. Subtracting the interval for
the anchor from that for the changed regime gives a valid optimum-response
interval. It is an enclosure in a specified model, not a confidence interval
for an estimated economic population. Strong competitors are evaluated by
the same construction.

\subsection{Trainable continuation factors}
For the square-activation critic, write
$L(W;P)=\frac14\|W'W-dP\|_F^2$, with symmetric own-policy target $P$.
Differentiation gives $\nabla_W L=W(W'W-dP)$. If $X\sim N(0,I)$, the mean
squared error between the gradients of $\|WX\|^2/d$ and $X'PX$ equals
$4\|W'W-dP\|_F^2/d^2$. The target is obtained by evaluating the current
candidate policy; it is not the Riccati value matrix. The experiment trains
all hidden weights for a fixed number of gradient updates and then solves
the actor's full vector-action problem. A conventional quadratic receives
the identical population target and is allowed its exact minimizer. This
matched-information comparison isolates the cost of the neural factorization.
No advantage over conventional quadratics is built into the target oracle.

\section{Recursive-risk and execution envelopes}
\label{sec:r27riskproof}
\begin{lemma}[Gaussian continuation perturbations]\label{lem:r27riskderivative}
Let $P\succeq L\succeq0$, $\|P-L\|_2\leq d_0$, and
$\gamma=1-2\theta\|\Sigma\|_2\|P\|_2>0$. Then
\[
 0\preceq\Psi_\theta(P)-\Psi_\theta(L)
       \preceq\frac{d_0}{\gamma^2}I,\qquad
 0\leq h_\theta(P)-h_\theta(L)
       \leq\frac{d_0\tr\Sigma}{\gamma}.
\]
The same inequalities hold at $\theta=0$ by continuity.
\end{lemma}
\begin{proof}
Gaussian integration after completing the square gives
\eqref{eq:r27gaussian}. Its derivatives in a symmetric direction $E$ are
\[
 D\Psi_\theta(P)[E]
 =(I-2\theta P\Sigma)^{-1}E(I-2\theta\Sigma P)^{-1},
\]
and
\[
 Dh_\theta(P)[E]
 =\tr\{\Sigma^{1/2}(I-2\theta\Sigma^{1/2}P\Sigma^{1/2})^{-1}
                      \Sigma^{1/2}E\}.
\]
Along the segment from $L$ to $P$, the inverse norms in the first display are
at most $\gamma^{-1}$ by the Neumann series. The symmetric inverse in the
second display is bounded by $\gamma^{-1}I$. Integrating the derivatives,
using $0\preceq P-L\preceq d_0I$, proves the claims. This proof also applies
to singular positive semidefinite $\Sigma$.
\end{proof}

\begin{proof}[Proof of Theorem~\ref{thm:r27entropic}]
Set $L_t(x)=x'(P_t-\delta_tI)x+c_t-\kappa_t$. The stated checks ensure that
every continuation coefficient used in this construction is positive
semidefinite and inside the Gaussian moment domain. Terminally $L_T=U_T$.
At date $t$, write $P=P_{t+1}$ and $L=P-\delta_{t+1}I$.
Let $S(P)$ denote the coefficient of the minimized current-action objective
with continuation coefficient $P$. Completion of squares gives
\[
 P_t-S(P)=D_t'H_t^{-1}D_t\preceq\zeta_t^2I/r_t.
\]
Let $K^L$ minimize the objective using $\Psi_\theta(L)$, and
$H^L=R_t+\beta B_t'\Psi_\theta(L)B_t$. The exact identity
\[
 H^L(K_t-K^L)
 =D_t+\beta B_t'\{\Psi_\theta(P)-\Psi_\theta(L)\}F_t
\]
and Lemma~\ref{lem:r27riskderivative} imply
$\|A_t-B_tK^L\|_2\leq\bar f_t$.
Evaluating the larger-continuation objective at $K^L$ therefore gives
\[
 0\preceq S(P)-S(L)\preceq\beta\bar f_t^2w_t I.
\]
Consequently $P_t-\delta_tI\preceq S(L)$. Own-policy constants satisfy
$c_t=\beta\{c_{t+1}+h_\theta(P_{t+1})\}$.
The constant part of \eqref{eq:r27riskenvelope} and the lemma give
\[
 c_t-\kappa_t\leq
 \beta\{c_{t+1}-\kappa_{t+1}+h_\theta(L)\}.
\]
Together these are precisely the Bellman-subsolution inequality
$L_t\leq\inf_u\{\ell_t+\beta\mathcal R_\theta L_{t+1}\}$.
Order preservation and backward induction imply $L_0\leq J_0^*$, including
comparison with arbitrary history-adapted controls of finite recursive cost.
The feasible stored policy gives $J_0^*\leq U_0$. Subtracting proves
\eqref{eq:r27riskgap}.
\end{proof}

\begin{proposition}[Execution errors under recursive preferences]
\label{prop:r27riskexecution}
Suppose the implemented action differs from $-K_tx$ by an adapted error
$e_t$ satisfying $\|e_t\|_2\leq\varepsilon_{\rm ex}\|x\|_2$ at every date.
Let $k_t\geq\|K_t\|_2$ and $s_t\geq\|R_t\|_2$.
Starting with $\delta_T^e=\kappa_T^e=0$, define
\begin{align*}
 \gamma_t^e&=1-2\theta\sigma(p_{t+1}+\delta_{t+1}^e),&
 \nu_t&=(p_{t+1}+\delta_{t+1}^e)/\gamma_t^e,\\
 a_t^e&=2\varepsilon_{\rm ex}k_ts_t+
             \varepsilon_{\rm ex}^2s_t
       +\beta\nu_t(2\varepsilon_{\rm ex}f_tb_t+
                         \varepsilon_{\rm ex}^2b_t^2),\\
 \delta_t^e&=\beta f_t^2\delta_{t+1}^e/(\gamma_t^e)^2+a_t^e,&
 \kappa_t^e&=\beta\{\kappa_{t+1}^e+
                    \tau\delta_{t+1}^e/\gamma_t^e\}.
\end{align*}
If every $\gamma_t^e>0$, the implemented policy's cost is at most
$U_0(x)+\delta_0^e\|x\|^2+\kappa_0^e$. Its excess cost over $J^*$ is thus
bounded by $(\delta_0+\delta_0^e)\|x\|^2+\kappa_0+\kappa_0^e$.
\end{proposition}
\begin{proof}
Assume the next implemented value is bounded by
$x'(P_{t+1}+\delta_{t+1}^eI)x+c_{t+1}+\kappa_{t+1}^e$.
Apply Lemma~\ref{lem:r27riskderivative} between $P_{t+1}$ and
$P_{t+1}+\delta_{t+1}^eI$, using the margin $\gamma_t^e$.
The unchanged transition contributes at most
$\beta f_t^2\delta_{t+1}^e\|x\|^2/(\gamma_t^e)^2$.
The action perturbation contributes at most
$(2\varepsilon_{\rm ex}k_ts_t+\varepsilon_{\rm ex}^2s_t)\|x\|^2$
to current cost and at most
$\beta\nu_t(2\varepsilon_{\rm ex}f_tb_t+
\varepsilon_{\rm ex}^2b_t^2)\|x\|^2$ to the continuation quadratic.
The log-determinant and constant contributions are bounded by
$\kappa_t^e$. Backward induction proves the upper envelope. The last assertion
uses Theorem~\ref{thm:r27entropic}. The action-error premise is an explicit
mathematical contract; it is not inferred merely from using floating-point
hardware on unbounded inputs.
\end{proof}

\paragraph{Verified numerical inverses.}
For an interval matrix $S$, propose a numerical inverse $C$ and bound
$r\geq\|I-SC\|_\infty$. If $r<1$, then
\[
 \|S^{-1}-C\|_\infty\leq\|C\|_\infty r/(1-r).
\]
This follows from $S^{-1}=C(I-(I-SC))^{-1}$. A matrix ball with that uniform
entry radius encloses the inverse. The implementation uses this test for
$S=I-2\theta\Sigma P$, together with the moment-domain test, outward scalar
products, and independently enclosed own-policy matrices. Failed inverse,
positivity, domain, or final tolerance checks remain failures to certify.

\paragraph{Certified current-action responses.}
In this convex capital economy the optimal current-action objective is
strongly convex with quadratic modulus at least $r_0$.
If a policy has a full recursive-cost gap at most $E$ at a given initial state,
its current action differs from the optimal action by at most $\sqrt{E/r_0}$.
The mean of its $d$ action coordinates therefore differs by at most
$\sqrt{E/(dr_0)}$. Pairing two such intervals gives a valid current-action
response band under a change in risk preferences with all future policies
reoptimized. No independence or population sampling assertion is needed.

\section{Proofs for Primitive-Budget Recursive Learning}\label{sec:r31proofs}
\subsection{The Gaussian continuation map}
For symmetric $P$ in the norm domain $2\theta\sigma\|P\|_2<1$, the map
$\Psi(P)=P(I-2\theta\Sigma P)^{-1}$ is symmetric. This follows from its
convergent series, or from
$(I-2\theta P\Sigma)^{-1}P=P(I-2\theta\Sigma P)^{-1}$;
commutation of $P$ and $\Sigma$ is not required. Differentiation gives
\[
 D\Psi(P)[E]=(I-2\theta P\Sigma)^{-1}
               E(I-2\theta\Sigma P)^{-1}.
\]
On a segment with $\|P\|_2\leq p$ and
$\gamma=1-2\theta\sigma p>0$, the Neumann bound therefore gives
$\|D\Psi(P)[E]\|_2\leq\gamma^{-2}\|E\|_2$.
The corresponding scalar term in the certainty equivalent has derivative
\[
 Dh_\theta(P)[E]=\operatorname{tr}
       \{(I-2\theta\Sigma P)^{-1}\Sigma E\}.
\]
For a positive-semidefinite segment, its norm is bounded by
$\operatorname{tr}(\Sigma)\|E\|_2/\gamma$. The limits at $\theta=0$ are
$D\Psi(P)[E]=E$ and $Dh_0(P)[E]=\operatorname{tr}(\Sigma E)$.
These distinct powers of $\gamma$ explain the separate coefficient and
constant weights in the policy certificate.

\subsection{Proof of Proposition~\ref{prop:r31tube}}
The assertion is exact at $T$. Suppose it holds at $t+1$. Both the proposed
center and the unknown exact coefficient lie in the norm ball of radius
$p_{t+1}$ about zero. The positive margin guarantees that the exact Gaussian
moment exists. The derivative bound on the segment joining the two matrices
implies
\[
 \|\Psi(P_{t+1})-\Psi(C_{t+1})\|_2
                         \leq e_{t+1}/\gamma_t^2.
\]
Subtract the proposed Bellman equation from exact own-policy evaluation.
The triangle inequality gives \eqref{eq:r31tube}. Positive definiteness of
the exact coefficient follows inductively from the positive stage cost and
the positive-semidefinite transformed continuation. The constants in the
own-policy values are finite because every Gaussian domain margin is
positive. No positivity of the numerical proposal itself is needed.

Writing the exact action residual as
$D_t=R_tK_t-\beta B_t'\Psi(P_{t+1})F_t$ and adding and subtracting the
centered expression proves \eqref{eq:r31defect}, using
$\|B'EF\|_F\leq\|B\|_2\|E\|_2\|F\|_F$.
All inequalities concern exact objects enclosed by the reported arithmetic;
the nominal inverse or nominal residual is never used without its enclosure.

\subsection{A trainable-factor contraction}
Let $mI\preceq M\preceq LI$, $W_0=\sqrt c\,O$ with $O$ orthogonal,
$0<c\leq m$, and $0<\alpha\leq1/(4L)$. Set $X_j=W_j'W_j$.
The iteration in \eqref{eq:r31training} implies
\[
 X_{j+1}=(I+\alpha(M-X_j))X_j(I+\alpha(M-X_j)).
\]
Since $X_0=cI$, every $X_j$ commutes with $M$. In their common eigenbasis,
for an eigenvalue $\lambda$ of $M$,
\[
 x_{j+1}=x_j\{1+\alpha(\lambda-x_j)\}^2.
\]
If $c\leq x_j\leq\lambda$, the update is no smaller than $x_j$ and no
larger than $\lambda$. To prove the latter statement, factor
\[
 \sqrt\lambda-\sqrt{x_j}\{1+\alpha(\lambda-x_j)\}
 =(\sqrt\lambda-\sqrt{x_j})
       \{1-\alpha\sqrt{x_j}(\sqrt\lambda+\sqrt{x_j})\}\geq0.
\]
The last factor is at least $1-2\alpha L>0$. Subtracting the scalar update
also yields
\[
 0\leq\lambda-x_{j+1}
  =(\lambda-x_j)
     \{1-2\alpha x_j-\alpha^2x_j(\lambda-x_j)\}
  \leq(1-\alpha c)(\lambda-x_j).
\]
Consequently $cI\preceq X_j\preceq M$ and
$\|X_j-M\|_F\leq(1-\alpha c)^j\sqrt d\,L$.
This proves the sufficient cap in \eqref{eq:r31cap}. The proof trains all
factor entries; simultaneous diagonalization is an argument about the
iterates, not an eigenvalue oracle supplied to the algorithm.

\subsection{Proof of Theorem~\ref{thm:r31constructive}}
We first prove the primitive envelope and target bounds jointly, backward
in time. At $T$, the exact coefficient is $Q_f$. Suppose
$0\preceq P_{t+1}\preceq\bar p_{t+1}I$ and, at nonterminal dates,
$P_{t+1}\succeq Q_{t+1}$. The risk-domain test then implies
\[
 dq_{t+1}I\preceq d\Psi(P_{t+1})
       \preceq d\bar p_{t+1}I/\bar\gamma_t
\]
for every fitted date. The contraction just proved gives a positive
semidefinite fitted matrix $\widehat P= W'W/d\preceq\Psi(P_{t+1})$ with
$\|\widehat P-\Psi(P_{t+1})\|_F\leq a_t/d$.
The fitted actor satisfies
\[
 K_t=(R_t+\beta B_t'\widehat P B_t)^{-1}
                         \beta B_t'\widehat P A_t,
 \qquad \|A_t-B_tK_t\|_2\leq f_t.
\]
Its exact residual therefore obeys
\[
 \|D_t\|_F
 =\|\beta B_t'(\widehat P-\Psi(P_{t+1}))F_t\|_F
 \leq\beta b_tf_ta_t/d\leq\sqrt{s r_tq_t}.
\]
At the terminal actor this residual is zero. Completing the square in the
true one-step objective shows that the fitted actor's loss against its true
greedy improvement is bounded by $s q_t\|x\|^2$. Comparing that minimum
with the fixed reference action $-L_tx$ gives the matrix inequality
\[
 P_t\preceq Q_t+L_t'R_tL_t+
         \beta(A_t-B_tL_t)'\Psi(P_{t+1})(A_t-B_tL_t)+s q_tI
       \preceq\bar p_tI.
\]
Here $s q_tI\preceq\eta Q_t$. The lower bound $P_t\succeq Q_t$ follows
from own-policy evaluation. This closes the induction without assuming a
bound on a target not yet generated.

It remains to bound the recursive policy loss. Apply the lower-subsolution
recursion \eqref{eq:r27riskenvelope}. Suppose inductively that
$\delta_{t+1}\leq s\lambda_{t+1}\leq q_*$ and
$\kappa_{t+1}\leq s\mu_{t+1}$. Its continuation margin is at least
$\bar\gamma_t$, its full-action residual is at most
$\sqrt{s r_tq_t}\leq\sqrt{\eta r_tq_t}$, and its lower-greedy closed-loop
bound is at most $\widetilde f_t$. Thus
\[
 \delta_t\leq s q_t\mathbf1_{\{t<T-1\}}
       +\frac{\beta\widetilde f_t^2}{\bar\gamma_t^2}s\lambda_{t+1}
       =s\lambda_t,
 \qquad \kappa_t\leq s\mu_t.
\]
The condition $s\lambda_t\leq q_*$ ensures positivity of every required
lower continuation because $P_t\succeq Q_t\succeq q_tI$.
The terminal allowance is zero, so the induction starts. The full adapted
policy comparison in Theorem~\ref{thm:r27entropic}, followed by
\eqref{eq:r31allocate}, proves \eqref{eq:r31constructivegap}.

\subsection{Numerical acceptance and work}
The finite cap bounds ideal hidden-weight updates, not processor time.
The implementation may use a smaller step, an enlarged cap, and periodically
spaced residual checks. A failed numerical training threshold or a failed
risk-domain test remains a failed attempt. For a returned rounded gain, the
outward residual-centered evaluation and the separately bounded action error
are the acceptance criteria. The experiment charges primitive-envelope
construction, own-policy transforms, all factor updates, actor solves,
residual checks, successful and failed verification, and durable output.
A structural comparator receives the same analytic Gaussian primitives and
is timed and verified separately. Access to those primitives is not evidence
of a sampling advantage in a model whose continuation must be simulated.

\section{From Training Geometry to Policy Accuracy}\label{sec:r32geometry}
The constructive schedule above controls a particular initialization. A returned
factor need not have been trained from that initialization. We therefore give
a second, a posteriori route which uses both first- and second-order information.
It rules out a singular stationary factor without assuming that its rank has
already been certified. Throughout this section $W\in\R^{d\times d}$ is the
trainable hidden-weight matrix, not a fixed dictionary.

For a symmetric target $M\succeq mI\succ0$, put
\[
 f_M(W)=\tfrac14\|W'W-M\|_F^2,\qquad G_M(W)=W(W'W-M).
\]
The Hessian is a self-adjoint operator on matrices equipped with the Frobenius
inner product. Its lower bound below must hold in every matrix direction;
a finite selection of curvature probes does not establish that premise.

\begin{theorem}[Gradient--curvature certificate]\label{thm:r32geometry}
Suppose $\nabla^2 f_M(W)\succeq-\kappa I$, where $0\leq\kappa<m$.
Then
\begin{equation}\label{eq:r32geometry}
 \sigma_{\min}(W)^2\geq\frac{m-\kappa}{3},\qquad
 \|W'W-M\|_F\leq
       \sqrt{\frac{3}{m-\kappa}}\,\|G_M(W)\|_F.
\end{equation}
In particular, a zero-gradient factor satisfying this curvature condition
represents the target exactly.
\end{theorem}
\begin{proof}
For $E=W'W-M$ and an arbitrary direction $Z$, differentiation gives
\[
 D^2f_M(W)[Z,Z]
   =\tfrac12\|W'Z+Z'W\|_F^2+\operatorname{tr}(E Z'Z).
\]
Let $s=\sigma_{\min}(W)$ and choose a corresponding pair of unit singular
vectors $u,v$, so that $Wv=su$ and $W'u=sv$. These relations also hold at
$s=0$, with left and right null vectors. Set $Z=uv'$. Then $\|Z\|_F=1$ and
\[
 -\kappa\leq D^2f_M(W)[Z,Z]=3s^2-v'Mv\leq3s^2-m.
\]
This proves the lower singular-value bound. Since $m>\kappa$, $W$ is
invertible, and $E=W^{-1}G_M(W)$. The Frobenius norm inequality proves the
second assertion.
\end{proof}

At $W=0$, the gradient vanishes for every target, although its approximation
error is generally large. The Hessian has a direction with curvature at most
$-m$. Thus the theorem excludes exactly the false implication from a small
gradient alone that would arise at this stationary point. It does not assert
that arbitrary nonconvex training attains the required second-order condition.
The finite construction in Theorem~\ref{thm:r31constructive} and the
certificate here are distinct results.

\subsection{An inexact own-policy target}
The policy evaluator may supply a symmetric approximation $\widetilde M$ rather
than the exact target. Suppose an independent evaluation account establishes
$\|M-\widetilde M\|_2\leq e$. Let
\[
 \|G_{\widetilde M}(W)\|_F\leq g,\qquad
 \nabla^2 f_{\widetilde M}(W)\succeq-\widetilde\kappa I,
 \qquad \|W\|_F\leq w.
\]

\begin{corollary}[Evaluation-aware factor error]\label{cor:r32inexact}
If $M\succeq mI$ and $\widetilde\kappa+e<m$, then
\begin{equation}\label{eq:r32inexact}
 \|W'W-M\|_F\leq
 (g+we)\sqrt{\frac{3}{m-\widetilde\kappa-e}}.
\end{equation}
A sufficient check for the target floor is
$\widetilde M\succeq(m+e)I$.
\end{corollary}
\begin{proof}
The gradient difference is $-W(M-\widetilde M)$, with Frobenius norm at most
$\|W\|_F e$. The difference of the Hessian operators maps $Z$ to
$-Z(M-\widetilde M)$ and has operator norm at most $e$. Hence the true
Hessian is bounded below by $-(\widetilde\kappa+e)I$ and the true gradient
norm by $g+we$. Apply Theorem~\ref{thm:r32geometry}. The sufficient target
floor follows from the spectral error bound. Notice that a spectral error
for the target multiplies the \emph{Frobenius} norm of $W$ here; replacing
it by its spectral norm without a dimension allowance would be invalid.
\end{proof}

For the recursive economy the exact target at date $t$ is
$M_t=d\Psi_\theta(P_{t+1}^{\pi})$, evaluated under the returned policy's
own future. Proposition~\ref{prop:r31tube} supplies the sufficient radius
$e_t^{M}=d e_{t+1}/\gamma_t^2$ around
$\widetilde M_t=d\Psi_\theta(C_{t+1})$; the approximation of that transformed
center, when numerical, contributes its own additional radius.
Neither a target inherited from a different policy nor a nominal matrix
inverse supplies this radius without its separate evaluation bound.

Write $a_t$ for the right side of \eqref{eq:r32inexact} divided by $d$,
$\widehat P_t=W_t'W_t/d$, and let
\[
 Z_t=R_tK_t-\beta B_t'\widehat P_t(A_t-B_tK_t)
\]
be the actor residual against the fitted continuation. The true full-action
residual then satisfies
\begin{equation}\label{eq:r32actor}
 \|D_t\|_F\leq\|Z_t\|_F+
   \beta\|B_t\|_2\|A_t-B_tK_t\|_2 a_t.
\end{equation}
Indeed, subtracting the two residuals leaves
$\beta B_t'(\widehat P_t-\Psi_\theta(P_{t+1}^{\pi}))
(A_t-B_tK_t)$. This identity proves the bound. Substitution in
Theorem~\ref{thm:r27entropic} gives a full adapted-policy certificate whenever
its domain, positivity, and implementation conditions also hold. The curvature
check is consequently a training diagnostic with an explicit route to economic
accuracy, not a replacement for the independent economic stopping rule.

\subsection{The terminal boundary and the iteration account}
The terminal coefficient is known, not learned. Under recursive risk
preferences, the last actor must nevertheless transform that known coefficient:
\begin{equation}\label{eq:r32terminal}
 K_{T-1}=\{R_{T-1}+\beta B_{T-1}'\Psi_\theta(Q_f)B_{T-1}\}^{-1}
              \beta B_{T-1}'\Psi_\theta(Q_f)A_{T-1}.
\end{equation}
This is the meaning of direct terminal evaluation. Substituting $Q_f$ for
$\Psi_\theta(Q_f)$ in this formula is justified at $\theta=0$, or when the
transform actually equals $Q_f$, but not in general. The scalar constant
$h_\theta(Q_f)$ does not change this actor. The frozen implementation already
uses \eqref{eq:r32terminal}; the current manuscript makes the boundary explicit.

The isotropic factor construction also admits the sharper contraction
\begin{gather*}
 \|W_j'W_j-M\|_F\leq(1-2\alpha c)^j\sqrt d\,L,\\
 0<c\leq\lambda_{\min}(M),\qquad
 0<\alpha\leq(4L)^{-1},\qquad\|M\|_2\leq L.
\end{gather*}
To see this, use the common eigenbasis in the constructive proof. With
$c\leq x_j\leq\lambda$,
\[
 \lambda-x_{j+1}=(\lambda-x_j)
   \{1-2\alpha x_j-\alpha^2x_j(\lambda-x_j)\}
 \leq(1-2\alpha c)(\lambda-x_j).
\]
Nonnegativity follows from the previously proved upper bound
$x_{j+1}\leq\lambda$, and $1-2\alpha c\geq1/2$. This sharper account
justifies a cap based on $-\log(1-2\alpha c)$ as well as the looser published
cap. It changes neither the frozen experimental budgets nor their interpretation
in rounded arithmetic.

For small matrices the replication code represents stored binary inputs as
exact rational numbers and checks the whole Hessian by a rational
positive-semidefiniteness elimination. Its certificates are conditional on the
declared target-radius bound, which must come from evaluation rather than from
this algebraic check. This implementation validates the new implication and
its failure cases. It is not offered as a scalable replacement for the
fifty-dimensional residual verifier, and its timing is not a primary economic
comparison.

\section{Certified Warm Starts under Changing Continuations}\label{sec:r33warm}
A continuation fitted to one economic future can be useful for another without
being accurate enough to reuse unchanged. The relevant distinction is between
an admissible initial factor and an admissible final policy. This section gives
a finite training account for the former and retains the independent economic
certificate for the latter. In contrast to the isotropic construction in
Section~\ref{sec:r31constructive}, the starting Gram matrix need not commute
with the new target. The result also allows explicitly bounded errors in each
hidden-weight update.

Let $M$ be a symmetric continuation target satisfying
$mI\preceq M\preceq LI$, where $0<m\leq L$. The target is fixed during
one fit, after the candidate's future policy has been finalized. A factor
$W\in\R^{p\times d}$, with $p\geq d$, represents the scalar continuation
$\|Wx\|^2/d$; every entry of $W$ is trainable. Write
$H_j=W_j'W_j$ and consider the implemented update
\begin{equation}\label{eq:r33update}
 W_{j+1}=W_j\{I-\alpha(H_j-M)\}+\Delta_j,
 \qquad \|\Delta_j\|_F\leq\nu.
\end{equation}
The perturbation account must cover the whole computed step, including
products and subtraction, not only the final storage conversion. The theorem
is a deterministic implication for all errors satisfying this account.

\begin{theorem}[Perturbation-stable warm-start training]\label{thm:r33warm}
Choose $0<r\leq m/2$ and $0<\alpha\leq\{2(L+r)\}^{-1}$, and set
\begin{align}
 b&=L+r, & q&=1-3\alpha m/4,\notag\\
 z&=2\sqrt b(1+\alpha r)\nu+\nu^2,&
 f&=z/(1-q).\label{eq:r33noise}
\end{align}
Suppose $\|H_0-M\|_F\leq e_0\leq r$ and $f\leq r$. Then every iterate
in \eqref{eq:r33update} satisfies
\begin{equation}\label{eq:r33envelope}
 \|H_j-M\|_F\leq f+q^j(e_0-f)\leq r.
\end{equation}
In particular, if $f<a<e_0$, it suffices to take
\begin{equation}\label{eq:r33cap}
 k=\left\lceil\frac{\log\{(e_0-f)/(a-f)\}}{-\log q}\right\rceil
\end{equation}
updates to achieve Gram error at most $a$. If $e_0\leq a$, no update is
required. No simultaneous diagonalization of $H_0$ and $M$ is assumed.
\end{theorem}

The proof in Section~\ref{sec:r33proofs} uses a Frobenius-norm contraction
for a Sylvester operator, rather than the common eigenbasis used by the cold
construction. It therefore covers rotated and overparameterized factors and
warm starts whose Gram matrices lie on either side of the target. The positive
noise floor is economically relevant: a requested continuation tolerance at
or below that floor is not certified by the finite-cap formula. More accurate
arithmetic or a different verified construction is then required; increasing
the iteration count alone does not justify acceptance.

\subsection{Transport and own-policy evaluation}
Suppose the old factor is accompanied by
$\|W_0'W_0-M_{\rm old}\|_F\leq e_{\rm old}$ and the newly evaluated target
by $\|M-M_{\rm old}\|_F\leq h$. The computable gate
\begin{equation}\label{eq:r33gate}
 e_{\rm old}+h\leq r
\end{equation}
certifies an admissible warm start with $e_0=e_{\rm old}+h$. A direct enclosure
of $\|W_0'W_0-M\|_F$ is an alternative, potentially sharper gate. Neither gate
certifies the old action as accurate in the changed economy. If the gate
fails, the controlled cold construction remains available; the failed gate
and all subsequent work are charged. A target obtained from the previous
regime's policy cannot replace the new policy's own continuation without a
separate transport bound.

A numerical evaluator may return a symmetric center $\widetilde M$ and an
independently verified spectral radius $e^M$ for the exact own-policy target
$M^\pi$. Apply Theorem~\ref{thm:r33warm} to $\widetilde M$, using certified
bounds on that center. If its final Gram error is at most $a$, then
\begin{equation}\label{eq:r33eval}
 \|W'W-M^\pi\|_2\leq a+e^M.
\end{equation}
In the recursive capital economy, $M_t^\pi=d\Psi_\theta(P_{t+1}^\pi)$.
The residual-centered evaluation in Proposition~\ref{prop:r31tube} supplies
$e_t^M=d e_{t+1}/\gamma_t^2$, with an additional allowance for evaluating the
transformed center when needed. This is a statement about the returned
policy, not the policy used to produce an old warm start.

For a returned actor $K_t$, put $F_t=A_t-B_tK_t$,
$\widehat P_t=W_t'W_t/d$, and
$Z_t=R_tK_t-\beta B_t'\widehat P_tF_t$. Its exact full-action residual obeys
\begin{equation}\label{eq:r33policy}
 \|D_t\|_F\leq\|Z_t\|_F+
 \frac{\beta\|B_t\|_2\|F_t\|_F}{d}(a_t+e_t^M).
\end{equation}
The Frobenius norm of $F_t$, rather than its spectral norm without a dimension
allowance, is essential when $e_t^M$ is only a spectral radius.
Substituting these residuals in Theorem~\ref{thm:r27entropic}, with its moment,
positivity, and implementation conditions, gives the same full adapted-policy
comparison as before. The training theorem changes the construction of the
candidate, not the economic target or the acceptance rule.

\subsection{Reliability and complete work}
The reliability statement is uniform over the class of initial factors and
update perturbations just specified. It is not a confidence interval inferred
from two or three training seeds. A randomized implementation supported
entirely in this verified class inherits the deterministic cap. Whether a
rounded evaluator, update routine, and returned policy satisfy the premises
is checked separately. Failed domain, transport, precision, or final-policy
checks remain failures to certify.

For a sequence of changed futures, the warm-start service account is
\begin{equation}\label{eq:r33work}
 C_{\rm warm}=C_{\rm initial}+\sum_\ell
 \{C_{\rm evaluate,\ell}+C_{\rm gate,\ell}
   +C_{\rm fit,\ell}+C_{\rm actor,\ell}
   +C_{\rm verify,\ell}+C_{\rm write,\ell}\}.
\end{equation}
Failed gates, cold fallbacks, and failed final checks enter this same sum.
The cap gives an upper bound on the number of factor updates in an admissible
warm fit. Comparing two upper bounds on training iterations does not establish
a lower realized cost than a competitor that stops early. An economic work
advantage requires complete measured costs, or an upper bound for the warm
procedure below an independently justified lower bound for its comparator.
In particular, the analytic structural solution in the quadratic economy
remains a legitimate comparator. This distinction preserves the adverse
frontiers of the earlier experiments while giving a positive, checkable
criterion for when an old neural continuation can be retrained safely.

The additional validation exercises below concern this training implication,
including noncommuting starts, rectangular factors, bounded perturbations,
transport failure, and unattainable precision requests. They are algebraic
and implementation checks, not additional independent economic experiments
and not a replacement for the previously frozen services.

\section{Proofs for Certified Warm Starts}\label{sec:r33proofs}
\subsection{A noncommuting Gram contraction}
Let $E=H-M$ and suppose $\|E\|_F\leq r$. Then
$(m-r)I\preceq H\preceq(L+r)I=bI$. The ideal, unperturbed factor step has
Gram error
\begin{equation}\label{eq:r33gramidentity}
 E^+=E-\alpha(EH+HE)+\alpha^2EHE.
\end{equation}
Diagonalize $H$, solely for this argument, with eigenvalues
$h_i\in[m-r,b]$. The linear map
$E\mapsto E-\alpha(EH+HE)$ multiplies each matrix entry in that basis by
$1-\alpha(h_i+h_j)$. These multipliers lie between zero and
$1-2\alpha(m-r)$ because $2\alpha b\leq1$. Frobenius invariance and
$\|EHE\|_F\leq br\|E\|_F$ give
\[
 \|E^+\|_F\leq
 \{1-2\alpha(m-r)+\alpha^2br\}\|E\|_F.
\]
Since $\alpha b\leq1/2$ and $r\leq m/2$,
\[
 2(m-r)-\alpha br\geq2m-\tfrac52r\geq\tfrac34m.
\]
Thus $\|E^+\|_F\leq q\|E\|_F$, with $q$ as in
\eqref{eq:r33noise}. Notice that $0<q<1$ and that the diagonalizing basis
is that of the current $H$, not a common basis for $H$ and $M$. This is why
no commutation assumption has entered the proof.

\subsection{Perturbations and the finite cap}
Write $U=W(I-\alpha E)$. Since $\|W\|_2^2=\|H\|_2\leq b$,
$\|U\|_2\leq\sqrt b(1+\alpha r)$. Adding $\Delta$ to the factor changes
the Gram matrix by
\[
 U'\Delta+\Delta'U+\Delta'\Delta.
\]
Its Frobenius norm is at most
$2\|U\|_2\|\Delta\|_F+\|\Delta\|_F^2\leq z$.
This estimate is valid for every rectangular $p\times d$ factor.
Consequently, while the radius condition holds,
$\|H_{j+1}-M\|_F\leq q\|H_j-M\|_F+z$.
The scalar recursion with initial value $e_0$ has the exact solution
$f+q^j(e_0-f)$. Each value is a convex combination of $e_0$ and $f$, both
at most $r$. Induction therefore proves that the radius condition continues
to hold and establishes \eqref{eq:r33envelope}. If $f<a<e_0$, solving the
scalar inequality for $j$ yields \eqref{eq:r33cap}. If $e_0\leq a$, the
initial factor already meets the requested Gram threshold.

\subsection{Transport and economic residuals}
The triangle inequality proves \eqref{eq:r33gate}. For an inexact target,
$\|W'W-M^\pi\|_2\leq\|W'W-\widetilde M\|_F+
\|\widetilde M-M^\pi\|_2$, proving \eqref{eq:r33eval}.
Finally,
\[
 D_t-Z_t=\beta B_t'\{\widehat P_t-\Psi_\theta(P_{t+1}^\pi)\}F_t.
\]
The inequality $\|B'EF\|_F\leq\|B\|_2\|E\|_2\|F\|_F$ proves
\eqref{eq:r33policy}. All state and action quantifiers, Gaussian moment
margins, lower-subsolution positivity checks, and implementation allowances
are those of Theorem~\ref{thm:r27entropic}; none is inferred from the
training gate. In particular, a residual against an old regime's continuation
must first pay its own-policy evaluation discrepancy.

\subsection{Numerical interpretation of the error account}
For a routine that computes the ideal matrix expression exactly and then
rounds each entry to a grid of spacing $h$, nearest rounding gives
$\nu\leq h\sqrt{pd}/2$. A floating-point implementation that also rounds
matrix products must add enclosures for those operations before invoking
this bound. The exact-rational regression tests use the former, explicitly
specified routine. Their success does not validate an unenclosed BLAS
roundoff account. Similarly, the stored scalar cap is checked by exact
rational comparison rather than trusting a rounded logarithm at a ceiling
boundary. This verification concerns the cap arithmetic, not the economic
model's unknown evaluation radius.

\section{Proofs for Budgeted Warm-Start Policy Construction}
\label{sec:r34proof}
\subsection{The mixed-norm residual}
Let $E_t=W_t'W_t-\widetilde M_t$ and
$V_t=\widetilde M_t-dS_t$. Then
\[
 D_t-Z_t=\frac{\beta}{d}B_t'(E_t+V_t)F_t.
\]
The inequalities $\|B'EF\|_F\leq\|B\|_2\|E\|_F\|F\|_2$ and
$\|B'VF\|_F\leq\|B\|_2\|V\|_2\|F\|_F$, together with
$\|F_t\|_F\leq\sqrt d\|F_t\|_2$, prove \eqref{eq:r34mixed}.
The spectral discrepancy also gives
\[
 \|\widehat S_t-S_t\|_2\leq(a_t+e_t^M)/d\leq q_*.
\]
These statements do not require the center or the Gram error to commute
with the true coefficient.

\subsection{A joint backward induction}
The coefficient bound is exact at the terminal date. Suppose
$0\preceq P_{t+1}^\pi\preceq\bar p_{t+1}I$. The positive primitive margin
implies $0\preceq S_t\preceq u_tI$, hence
$0\preceq\widehat S_t\preceq(u_t+q_*)I$. The actor identity can be written
\[
 (R_t+\beta B_t'\widehat S_tB_t)K_t
                      =\beta B_t'\widehat S_tA_t+Z_t.
\]
Its left coefficient is bounded below by $r_tI$. Since
$z_t\leq\sqrt{s r_tq_t}\leq v_t$, it follows that
\[
 \|F_t\|_2\leq\|A_t\|_2+
    \frac{b_t}{r_t}\{\beta b_t(u_t+q_*)\|A_t\|_2+z_t\}
       \leq f_t^+.
\]
Equation~\eqref{eq:r34mixed} and the second local acceptance test therefore
yield $\|D_t\|_F\leq\sqrt{s r_tq_t}$. Completing the square in the exact
one-step action objective bounds the returned actor's excess coefficient
over its true greedy improvement by $s q_t I$. Comparison of that minimum
with the reference feedback $L_t$ gives
\[
 P_t^\pi\preceq Q_t+L_t'R_tL_t+
 \beta(A_t-B_tL_t)'S_t(A_t-B_tL_t)+s q_tI\preceq\bar p_tI,
\]
because $s q_tI\preceq\eta Q_t$. Positive stage costs give
$P_t^\pi\succeq Q_t\succeq q_tI$. This closes the coefficient induction.
In particular, the target bound used in the next fit is not inferred from
a presumed successful policy comparison.

Now apply the subsolution recursion of
Theorem~\ref{thm:r27entropic}. If
$\delta_{t+1}\leq s\lambda_{t+1}^+\leq q_*$, its true margin is at least
$\bar\gamma_t$ and its lower-greedy closed-loop bound is at most $g_t$.
Consequently
\[
 \delta_t\leq s q_t+
       \frac{\beta g_t^2}{\bar\gamma_t^2}s\lambda_{t+1}^+
           =s\lambda_t^+,
 \qquad
 \kappa_t\leq\beta\left(s\mu_{t+1}^+
       +\frac{\tau s\lambda_{t+1}^+}{\bar\gamma_t}\right)=s\mu_t^+.
\]
The zero terminal allowances start this induction. The allocation preserves
the lower subsolution's required positivity at every noninitial date.
The recursive full-policy theorem and the initial-radius inequality then
prove \eqref{eq:r34policybound}. Neither state sampling nor a restriction to
linear competitors enters this argument.

\subsection{Finite arithmetic and fallback}
For rational matrices and a finalized rational feedback, the Gaussian
coefficient map uses matrix multiplication and an inverse of
$I-2\theta\Sigma P$. A strict norm-domain margin implies invertibility.
The own-policy coefficient and the next fitted actor can therefore be
computed by finitely many rational operations. The state-independent
log-determinant constant is not an input to these actor equations; its
contribution to the policy comparison is already paid by $\mu_t^+$.
The scalar bounds can be enlarged rationally, and squared inequalities can
be checked without treating rounded square roots as exact.

With exact evaluation and actor solves, take $e_t^M=z_t=0$ and select a
positive rational $a_t$ satisfying \eqref{eq:r34accept}. The exact-product,
nearest-dyadic update has perturbation norm at most
$2^{-b}\sqrt{pd}/2$ at $b$ storage bits. Its Gram noise floor tends to zero
as $b$ tends to infinity. For a certified warm basin, choose $b$ so that the
floor is below $a_t$ and apply Theorem~\ref{thm:r33warm}. Its recurrence is
checked with rational upper bounds for the square roots. If the initial
error already lies below $a_t$, no factor update is needed.

On failure of the warm gate, choose rational $c=s_0^2$ with $0<c\leq m_t$
and $W_0=s_0I$ at a nonterminal date. The controlled cold contraction in
Section~\ref{sec:r31proofs}, evaluated without rounding its products, reaches
the same positive tolerance in finitely many rational updates. This is an
existence and termination statement for a specified arithmetic procedure,
not a practical bit-complexity bound for all exact rational implementations.
It does not justify discarding failed gates from a measured service cost.
The terminal transformed coefficient is supplied directly. These observations
prove Corollary~\ref{cor:r34finite}.

\section{Proofs for Precision-Adaptive Refresh}\label{sec:r37proof}
\begin{proof}[Proof of Theorem~\ref{thm:r37precision}]
Write $Y=WM^{-1/2}$ and $C=Y'Y=I+E$. The ideal update in
\eqref{eq:r37update}, followed by right multiplication by $M^{-1/2}$,
is $\overline Y=Y(3I-C)/2$. Consequently
\[
 \overline Y'\overline Y=C(3I-C)^2/4,
 \qquad \overline Y'\overline Y-I=(-3E^2+E^3)/4.
\]
These identities hold for rectangular $Y$; they require no commutation
between $W'W$ and $M$. The matrices in the right-hand identities commute
because they are polynomials in the single symmetric matrix $E$.
If $\|E\|_2\leq\rho\leq r$, the second identity gives
$\|\overline Y'\overline Y-I\|_2\leq c_r\rho^2$.
Moreover, for every eigenvalue $e\in[-r,r]$,
\[
 0\leq(1+e)(2-e)^2/4=1+e^2(e-3)/4\leq1.
\]
Hence $\|\overline Y\|_2\leq1$. Put $Z=\Delta M^{-1/2}$.
The implemented Gram matrix differs from the ideal one by
$\overline Y'Z+Z'\overline Y+Z'Z$, of norm at most $2\nu+\nu^2$.
The allocated precision therefore implies
\[
 \|E_{j+1}\|_2\leq c_r\rho_j^2+2\nu_j+\nu_j^2
                         \leq\rho_j^2=\rho_{j+1}.
\]
Induction proves \eqref{eq:r37double}, since $\rho_j^2\leq\rho_j\leq r$.
Solving $\rho_0^{2^K}\leq\tau$ gives \eqref{eq:r37cap}. The zero-update
cases are immediate.

Entrywise nearest-grid rounding has Frobenius error at most
$\sqrt{pd}\,2^{-b-1}$. Multiplication by $M^{-1/2}$ and
$\|M^{-1/2}\|_2\leq m^{-1/2}$ give the specified upper bound $\nu$.
In particular, $\nu\leq(1-c_r)\rho_j^2/3$ is sufficient: it implies
$\nu\leq1$ and thus $2\nu+\nu^2\leq3\nu\leq(1-c_r)\rho_j^2$.
An admissible integer therefore always exists, and the first admissible
positive integer satisfies
\[
 b_j\leq C+2\log_2(1/\rho_j)
\]
for a finite constant $C$ depending only on $A$ and $r$. If $K\geq1$ is
the first successful update count, $\rho_0^{2^{K-1}}>\tau$. Hence
\begin{align*}
 \sum_{j<K} b_j
 &\leq CK+2\log_2(1/\rho_0)(2^K-1)\\
 &<CK+4\log_2(1/\tau),
\end{align*}
which proves \eqref{eq:r37bitaccount}. A zero-update fit contributes an
empty sum. This bound is on stored fractional precision, not on the bit
lengths of intermediate rational products.

Finally, the target is fixed throughout the fit. Its exact inverse can be
constructed once, checked against $M$, and applied to each successive Gram
matrix. The inverse applications and their residual checks are still
performed and charged. No inverse is needed when the initial factor is
accepted. This establishes the final accounting assertion without claiming
that the training or the checks have zero cost.
\end{proof}

\begin{proof}[Proof of the target-change bound]
The old certificate and target comparison give
\[
 (1-\rho_{\rm old})(1-\omega)M_{\rm new}
 \preceq W'W\preceq
 (1+\rho_{\rm old})(1+\omega)M_{\rm new}.
\]
The upper relative excess is $\rho_{\rm old}+\omega+
\rho_{\rm old}\omega$; the lower relative deficit is
$\rho_{\rm old}+\omega-\rho_{\rm old}\omega$, which is no greater.
The larger quantity thus certifies the proposed new basin radius. The
calculation concerns only the old factor, not an inverse or a policy that
has been accepted under a different target.
\end{proof}

\begin{proof}[Proof of Corollary~\ref{cor:r37policy}]
For the current center $\widetilde M_t$, the relative Gram enclosure yields
\[
 \|W_t'W_t-\widetilde M_t\|_F
 \leq\sqrt d\,\|\widetilde M_t\|_2\|E_t\|_2
 \leq\sqrt d\,L_t\tau_t\leq a_t.
\]
The target-change gate, or a direct gate, supplies the initial basin. The
precision theorem supplies this inequality in finitely many steps for each
positive allowance, provided the implementation supplies the required
precision. The terminal transformed coefficient is evaluated directly.
Working backward after finalizing each future policy therefore supplies
exactly the Gram, evaluation, and actor premises of
Theorem~\ref{thm:r34robust}. That theorem proves the coefficient envelope
and the all-action Bellman comparison jointly, and gives the displayed
initial-state policy bound. No inference from sampled states, a finite
action catalogue, or a small fitting loss is used in this implication.
An independently bounded discrepancy between the stored policy and its
executed actions must be added through its own implementation account.
\end{proof}


\section{Complete factorial and timing ledger}\label{sec:r39ledger}
The following entries summarize all 96 quadratic economy--method objects. Each has three isolated repetitions. Times are complete-service seconds; the range gives the observed minimum and maximum. $U$ counts updates, $B$ target-inverse builds and $F$ failed policy checks in the final tuning trial. The machine-readable audit and original records additionally preserve all tuning trials, actor solves, own-policy evaluation dates, moment margins, native precision, process memory and startup-inclusive clocks. Methods: FR/FC are fixed64 rebuilt/cached; AR/AC are adaptive rebuilt/cached; TC is tuned native cached; S is structural.
\begin{longtable}{llrrrrrr}
\toprule Object & Method & Median & Min & Max & $U$ & $B$ & $F$\\\midrule\endhead
tol 0.0001 & AC & 0.433 & 0.417 & 0.433 & 11 & 3 & 15\\
tol 0.0001 & AR & 0.421 & 0.418 & 0.422 & 11 & 11 & 15\\
tol 0.0001 & FC & 0.415 & 0.415 & 0.421 & 11 & 3 & 15\\
tol 0.0001 & FR & 0.416 & 0.416 & 0.420 & 11 & 11 & 15\\
tol 0.0001 & S & 0.170 & 0.168 & 0.173 & 0 & 0 & 4\\
tol 0.0001 & TC & 0.424 & 0.423 & 0.428 & 11 & 3 & 15\\
tol 0.01 & AC & 0.352 & 0.350 & 0.356 & 8 & 3 & 12\\
tol 0.01 & AR & 0.351 & 0.349 & 0.358 & 8 & 8 & 12\\
tol 0.01 & FC & 0.353 & 0.352 & 0.354 & 8 & 3 & 12\\
tol 0.01 & FR & 0.354 & 0.352 & 0.354 & 8 & 8 & 12\\
tol 0.01 & S & 0.175 & 0.172 & 0.181 & 0 & 0 & 4\\
tol 0.01 & TC & 0.356 & 0.351 & 0.356 & 8 & 3 & 12\\
tol 1e-06 & AC & 0.444 & 0.443 & 0.450 & 12 & 3 & 16\\
tol 1e-06 & AR & 0.441 & 0.441 & 0.449 & 12 & 12 & 16\\
tol 1e-06 & FC & 0.442 & 0.440 & 0.444 & 12 & 3 & 16\\
tol 1e-06 & FR & 0.448 & 0.440 & 0.456 & 12 & 12 & 16\\
tol 1e-06 & S & 0.171 & 0.170 & 0.172 & 0 & 0 & 4\\
tol 1e-06 & TC & 0.451 & 0.438 & 0.453 & 12 & 3 & 16\\
tol 1e-08 & AC & 0.500 & 0.485 & 0.503 & 14 & 3 & 18\\
tol 1e-08 & AR & 0.489 & 0.483 & 0.500 & 14 & 14 & 18\\
tol 1e-08 & FC & 0.482 & 0.481 & 0.484 & 14 & 3 & 18\\
tol 1e-08 & FR & 0.487 & 0.485 & 0.495 & 14 & 14 & 18\\
tol 1e-08 & S & 0.169 & 0.169 & 0.170 & 0 & 0 & 4\\
tol 1e-08 & TC & 0.491 & 0.486 & 0.492 & 14 & 3 & 18\\
tol 1e-10 & AC & 0.486 & 0.483 & 0.499 & 14 & 3 & 18\\
tol 1e-10 & AR & 0.487 & 0.486 & 0.491 & 14 & 14 & 18\\
tol 1e-10 & FC & 0.491 & 0.484 & 0.509 & 14 & 3 & 18\\
tol 1e-10 & FR & 0.491 & 0.486 & 0.500 & 14 & 14 & 18\\
tol 1e-10 & S & 0.175 & 0.173 & 0.357 & 0 & 0 & 4\\
tol 1e-10 & TC & 0.488 & 0.488 & 0.502 & 14 & 3 & 18\\
moment & AC & 0.490 & 0.488 & 0.507 & 14 & 3 & 18\\
moment & AR & 0.492 & 0.490 & 0.501 & 14 & 14 & 18\\
moment & FC & 0.479 & 0.476 & 0.491 & 13 & 3 & 17\\
moment & FR & 0.474 & 0.469 & 0.479 & 13 & 13 & 17\\
moment & S & 0.171 & 0.171 & 0.178 & 0 & 0 & 4\\
moment & TC & 0.495 & 0.495 & 0.507 & 14 & 3 & 18\\
grid 16-8-16.0 & AC & 2.481 & 2.476 & 2.533 & 42 & 7 & 50\\
grid 16-8-16.0 & AR & 2.474 & 2.463 & 2.636 & 42 & 42 & 50\\
grid 16-8-16.0 & FC & 2.542 & 2.459 & 2.543 & 42 & 7 & 50\\
grid 16-8-16.0 & FR & 2.518 & 2.471 & 2.568 & 42 & 42 & 50\\
grid 16-8-16.0 & S & 0.476 & 0.470 & 0.479 & 0 & 0 & 8\\
grid 16-8-16.0 & TC & 2.493 & 2.456 & 2.541 & 42 & 7 & 50\\
grid 32-4-16.0 & AC & 0.859 & 0.856 & 0.876 & 18 & 3 & 22\\
grid 32-4-16.0 & AR & 0.855 & 0.850 & 0.875 & 18 & 18 & 22\\
grid 32-4-16.0 & FC & 0.905 & 0.887 & 0.909 & 19 & 3 & 23\\
grid 32-4-16.0 & FR & 0.892 & 0.889 & 0.895 & 19 & 19 & 23\\
grid 32-4-16.0 & S & 0.225 & 0.224 & 0.226 & 0 & 0 & 4\\
grid 32-4-16.0 & TC & 0.857 & 0.855 & 0.858 & 18 & 3 & 22\\
grid 4-12-4.0 & AC & 3.683 & 3.648 & 3.708 & 44 & 11 & 56\\
grid 4-12-4.0 & AR & 3.612 & 3.553 & 3.614 & 44 & 44 & 56\\
grid 4-12-4.0 & FC & 3.645 & 3.556 & 3.725 & 44 & 11 & 56\\
grid 4-12-4.0 & FR & 3.636 & 3.533 & 3.654 & 44 & 44 & 56\\
grid 4-12-4.0 & S & 0.856 & 0.839 & 0.861 & 0 & 0 & 12\\
grid 4-12-4.0 & TC & 3.633 & 3.551 & 3.752 & 44 & 11 & 56\\
grid 8-4-8.0 & AC & 0.532 & 0.527 & 0.549 & 15 & 3 & 19\\
grid 8-4-8.0 & AR & 0.536 & 0.534 & 0.540 & 15 & 15 & 19\\
grid 8-4-8.0 & FC & 0.534 & 0.526 & 0.538 & 15 & 3 & 19\\
grid 8-4-8.0 & FR & 0.527 & 0.526 & 0.534 & 15 & 15 & 19\\
grid 8-4-8.0 & S & 0.178 & 0.176 & 0.180 & 0 & 0 & 4\\
grid 8-4-8.0 & TC & 0.529 & 0.528 & 0.537 & 15 & 3 & 19\\
change 0.25-cold & AC & 0.448 & 0.446 & 0.449 & 12 & 3 & 16\\
change 0.25-cold & AR & 0.450 & 0.448 & 0.454 & 12 & 12 & 16\\
change 0.25-cold & FC & 0.445 & 0.445 & 0.445 & 12 & 3 & 16\\
change 0.25-cold & FR & 0.445 & 0.441 & 0.447 & 12 & 12 & 16\\
change 0.25-cold & S & 0.172 & 0.168 & 0.174 & 0 & 0 & 4\\
change 0.25-cold & TC & 0.446 & 0.437 & 0.451 & 12 & 3 & 16\\
change 0.25-warm & AC & 0.872 & 0.863 & 0.880 & 12 & 3 & 16\\
change 0.25-warm & AR & 0.864 & 0.858 & 0.901 & 12 & 12 & 16\\
change 0.25-warm & FC & 0.863 & 0.853 & 0.865 & 12 & 3 & 16\\
change 0.25-warm & FR & 0.882 & 0.876 & 0.919 & 12 & 12 & 16\\
change 0.25-warm & S & 0.598 & 0.584 & 0.602 & 0 & 0 & 4\\
change 0.25-warm & TC & 0.872 & 0.869 & 0.884 & 12 & 3 & 16\\
change 1.005-cold & AC & 0.439 & 0.437 & 0.442 & 12 & 3 & 16\\
change 1.005-cold & AR & 0.448 & 0.442 & 0.473 & 12 & 12 & 16\\
change 1.005-cold & FC & 0.441 & 0.436 & 0.460 & 12 & 3 & 16\\
change 1.005-cold & FR & 0.440 & 0.439 & 0.473 & 12 & 12 & 16\\
change 1.005-cold & S & 0.173 & 0.169 & 0.173 & 0 & 0 & 4\\
change 1.005-cold & TC & 0.447 & 0.441 & 0.464 & 12 & 3 & 16\\
change 1.005-warm & AC & 0.679 & 0.658 & 0.680 & 3 & 3 & 7\\
change 1.005-warm & AR & 0.682 & 0.668 & 0.691 & 3 & 3 & 7\\
change 1.005-warm & FC & 0.659 & 0.657 & 0.677 & 3 & 3 & 7\\
change 1.005-warm & FR & 0.665 & 0.662 & 0.689 & 3 & 3 & 7\\
change 1.005-warm & S & 0.605 & 0.602 & 0.622 & 0 & 0 & 4\\
change 1.005-warm & TC & 0.666 & 0.660 & 0.667 & 3 & 3 & 7\\
change 4.0-cold & AC & 0.451 & 0.440 & 0.452 & 12 & 3 & 16\\
change 4.0-cold & AR & 0.444 & 0.443 & 0.452 & 12 & 12 & 16\\
change 4.0-cold & FC & 0.446 & 0.436 & 0.456 & 12 & 3 & 16\\
change 4.0-cold & FR & 0.441 & 0.440 & 0.448 & 12 & 12 & 16\\
change 4.0-cold & S & 0.172 & 0.171 & 0.173 & 0 & 0 & 4\\
change 4.0-cold & TC & 0.450 & 0.437 & 0.462 & 12 & 3 & 16\\
change 4.0-warm & AC & 0.859 & 0.856 & 0.882 & 12 & 3 & 16\\
change 4.0-warm & AR & 0.873 & 0.866 & 0.876 & 12 & 12 & 16\\
change 4.0-warm & FC & 0.872 & 0.863 & 0.944 & 12 & 3 & 16\\
change 4.0-warm & FR & 0.871 & 0.870 & 0.906 & 12 & 12 & 16\\
change 4.0-warm & S & 0.594 & 0.588 & 0.604 & 0 & 0 & 4\\
change 4.0-warm & TC & 0.870 & 0.869 & 0.890 & 12 & 3 & 16\\
\bottomrule\end{longtable}
\subsection{Nonlinear timing dispersion}
\begin{tabular}{lrrrr}\toprule Object & Median & Min & Max & Peak RSS (KiB)\\\midrule
nonlinear-frontier & 6.077 & 5.993 & 6.086 & 45700\\
neural-1.0-0.0 & 2.846 & 2.840 & 2.848 & 313708\\
neural-1.0-1.0 & 5.011 & 4.984 & 6.237 & 314288\\
neural-4.0-0.0 & 2.999 & 2.976 & 3.005 & 313744\\
neural-4.0-1.0 & 6.222 & 6.220 & 6.225 & 314420\\
capital-1.0-0.0 & 0.240 & 0.240 & 0.241 & 45400\\
capital-1.0-1.0 & 1.311 & 1.310 & 1.325 & 45400\\
capital-4.0-0.0 & 0.409 & 0.408 & 0.410 & 45400\\
capital-4.0-1.0 & 2.564 & 2.562 & 2.569 & 45400\\
\bottomrule\end{tabular}
\par\medskip The frontier is one complete four-mesh service per repetition. Its internal cumulative construction clocks are distinct from this full-service clock. Peak RSS is process high-water memory, not exclusive algorithm storage.

\bibliographystyle{ecta-fullname}
\bibliography{revisions/2026-10-07-r39/supp_references}
\end{document}
